\documentclass[12pt,a4paper]{article}

% --- Pacotes Fundamentais ---
\usepackage[utf8]{inputenc}
\usepackage[T1]{fontenc}
\usepackage{lmodern}
\usepackage[brazil,english]{babel}
\usepackage{geometry}
\geometry{a4paper, margin=2.5cm}
\usepackage{amsmath,amssymb,amsfonts}
\usepackage{graphicx}
\usepackage{booktabs}
\usepackage{microtype}
\usepackage{url}

% --- Metadados do Manuscrito ---
\title{\textbf{Título da Pesquisa Científica: Subtítulo Explicativo}}
\author{
    \textbf{Vitor}\\
    Programa de Pós-Graduação / Pesquisa em TI\\
    \texttt{autor@exemplo.com}
}
\date{\today}

\begin{document}
\selectlanguage{brazil}

\maketitle

\begin{abstract}
Este é o resumo da pesquisa acadêmica. O resumo deve apresentar de forma concisa e objetiva o contexto do problema, a lacuna teórica ou prática identificada, o objetivo principal do trabalho, a metodologia adotada (qualitativa, quantitativa ou métodos mistos) e os principais resultados obtidos ou contribuições esperadas.
\end{abstract}

\textbf{Palavras-chave:} Pesquisa Científica; Metodologia; Engenharia de Software; Ciência da Computação.

\vspace{1em}
\hrule
\vspace{1.5em}

% --- Inclusão Modular de Seções ---
\section{Introdução}
\label{sec:introducao}

A formulação rigorosa de um problema de pesquisa é o primeiro passo para qualquer investigação acadêmica relevante \cite{gil2022projetos, creswell2022research}. Nesta seção inicial, delimita-se o contexto temático, o escopo da pesquisa e a justificativa para a condução do estudo.

\subsection{Contextualização e Motivação}
Na área de Tecnologia da Informação, o avanço veloz de ferramentas, modelos e arquiteturas exige abordagens metodológicas sólidas para distinguir inovações sustentáveis de tendências passageiras \cite{oates2022researching}.

\subsection{Problema de Pesquisa e Objetivos}
O objetivo geral deste trabalho consiste em analisar e estruturar um fluxo de trabalho reprodutível para apoiar a condução e publicação de pesquisas acadêmicas.

\section{Trabalhos Relacionados}
\label{sec:relacionados}

A revisão sistemática da literatura constitui a base epistemológica sobre a qual novas teorias ou experimentos são desenvolvidos \cite{booth2021systematic, boland2017systematic}. 

A arte da pesquisa requer o diálogo explícito entre argumentos consolidados e novas proposições acadêmicas \cite{booth2024craft, graff2024theysay}. Nesta seção, mapeiam-se os principais autores e métodos correlatos já reportados na literatura científica.

\section{Metodologia}
\label{sec:metodologia}

A metodologia adotada segue as diretrizes consagradas de delineamento de pesquisa \cite{marconi2017metodologia, creswell2022research}. Define-se aqui o tipo de abordagem (qualitativa, quantitativa ou mista), os procedimentos técnicos de coleta de dados e os critérios de validação.

\subsection{Delineamento Experimental}
Para experimentos computacionais e análise estatística, aplicam-se modelos formais de aprendizado e métricas padronizadas \cite{james2023statistical}.

\section{Resultados e Discussão}
\label{sec:resultados}

Nesta seção são reportados os achados empíricos e as análises interpretativas decorrentes da aplicação da metodologia.

\subsection{Apresentação dos Dados}
Recomenda-se a utilização de tabelas formatadas de acordo com as normas científicas (usando o pacote \texttt{booktabs}) e figuras nítidas e vetorizadas.

\section{Conclusão}
\label{sec:conclusao}

A conclusão retoma sinteticamente os objetivos delineados, avalia o atendimento às hipóteses ou questões de pesquisa, elenca as principais contribuições teóricas e práticas e aponta direções para trabalhos futuros.


% --- Referências Bibliográficas ---
\newpage
\bibliographystyle{plain}
\bibliography{references}

\end{document}
